% New English translation from checked cote 25, site 16 (PDF page 17).
% AI-translated by GPT-6 Astra (OpenAI, ChatGPT Pro mode).
\scanpage{25/16}
\mspage{\hand{13}}

\var{\deleted{\uline{Corollary} 2.8. Under the conditions of 2.7., for $Z_i$ to be
nonembedded, it is n and s that $Z_i$ \var{\inserted[typed]{over \unsure{$\bar{k}$}}}{over \unsure{$\bar{k}$}}{insertion}{ega25_25_016-v0002} be nonembedded. When this is so, the
radicial multiplicity of $F$ at $Z_i$ equals the radicial
multiplicity of $G$ over $k(\ \ )$ at $Z_i$.

The first assertion follows immediately from the last assertion
of 2.7. As for the first, we may assume\var{\xout{, (since this is a}}{}{typewriter strike-out}{ega25_25_016-v0003}
\var{\inserted[typed]{by restriction to an open subset of X and}}{by restriction to an open subset of X and}{insertion}{ega25_25_016-v0004}
\var{\xout{\unsure{geometric notion}}}{}{typewriter strike-out}{ega25_25_016-v0005} by dévissage on $F$ (since $Y$ is flat over $X$
and therefore $F\ \ G$ is an exact functor) that $F=\underline{O}_{Z_i}$, and we can
clearly replace $X$ by $Z_i$. We are thus reduced to the case where $X$ is
integral and $F=\underline{O}_X$.}}{}{deletion}{ega25_25_016-v0001} Let $K/k$ be the function field of $X$, and let $n=\dim P$;
introducing \var{\xout{projec}}{}{typewriter strike-out}{ega25_25_016-v0007} affine coordinates $T_1,..,T_n$ in $P$ (by choosing
a hyperplane at infinity $H_\hand{\infty}$ such that $f(X)$ is not contained in
it) and $S_1,...S_n$ in $\overset{\hand{\vee}}{P}$, we see that the field\var{\xout{\unsure{ex}}}{}{typewriter strike-out}{ega25_25_016-v0008}\var{\inserted[typed]{of}}{ of}{insertion}{ega25_25_016-v0009} functions $L$
of $Y_{\hand{\eta}}$ \var{\xout{xxx}}{}{typewriter strike-out}{ega25_25_016-v0010} identifies with the fraction field of the integral ring
$K[S_1,..,S_n]/(\hand{\sum} t_iS_i - 1)$, where the $t_i \hand{\in} K$ are the images of the $T_i$ by
means of $f:X\hand{\to} P$. Since $\dim \hand{\overline{\textcolor{black}{f(X)}}} \hand{>} 0$, the $t_i$ are not all
algebraic over $k$, and a fortiori they are not all zero; suppose,
e.g., $t_n \hand{\neq} 0$. We then immediately find $L = K(S_1,..,S_{n-1})$,
(a purely transcendental extension), $S_n\hand{\in} L$ being given by the equation
\[ \hand{\sum} t_iS_i - 1 = 0 \]
in terms of the $S_i$ $(1\hand{\leq} i\hand{\leq} n-1)$ and the $t_i$ $(1\hand{\leq} i\hand{\leq} n)$. On the other hand,
$k' = k(\hand{\eta})$ identifies with $k(S_1,..,S_n)$, and the canonical inclusion
$k'\hand{\longrightarrow} L$ is obtained by sending $S_i$ to $S_i$, i.e. $k'$, as a sub\var{\xout{\unsure{exten}}}{}{typewriter strike-out}{ega25_25_016-v0012}\var{\inserted[typed]{extension}}{extension}{insertion}{ega25_25_016-v0013} of $L$, is the subextension generated by the $S_i$ $(1\hand{\leq} i\hand{\leq} n)$, or,
what clearly amounts to the same thing, by the $S_i$ $(1\hand{\leq} i\hand{\leq} \var{\xout{n}}{}{typewriter strike-out}{ega25_25_016-v0014}\var{\inserted[typed]{n-1}}{n-1}{insertion}{ega25_25_016-v0015})$ and by
$S_n = a_o + a_1S_1 + \ldots + a_{n-1}S_{n-1}$, where $a_o = t_n^{-1}$, $a_i =-t_it_n^{-1}$ for
$1\hand{\leq} i\hand{\leq} n-1$. \var{\deleted{Since}}{}{deletion}{ega25_25_016-v0016}\var{\inserted{Note that}}{Note that}{insertion}{ega25_25_016-v0017} the field generated by the $a_i$ or by the $t_i$ is clearly
the same, \var{\deleted{the $a_i$ are likewise \unsure{not|not} all algebraic over k.}}{}{deletion}{ega25_25_016-v0018}\var{\inserted{their common transcendence degree
is precisely $\dim f(X)$.}}{their common transcendence degree
is precisely $\dim f(X)$.}{insertion}{ega25_25_016-v0019}
\tnote{25-016-1}{The cancelled proof says first a second time, and leaves a blank between F and G; both features are retained. The cancelled s after le and a inside l'inclusion do not change English the or inclusion and have no English variants. The doubtful ne/n in the cancelled sentence is represented by the repeated candidate not. The source splits choisiss ant; its l'adde is understood in the phrase by means of, with both spellings preserved in the French.}

\var{\margin{16}}{\placedmargin{16}}{marginal note placed}{ega25_25_016-v0020}
